\documentclass[aspectratio=169]{beamer}
\usetheme{metropolis}

\usepackage[utf8]{inputenc}
\usepackage[spanish]{babel}
\usepackage{graphicx}
\usepackage{xcolor}
\usepackage{tikz}
\usetikzlibrary{positioning, arrows.meta}
\usepackage{tcolorbox}
\tcbuselibrary{skins, breakable}
\usepackage{fontawesome5}

% ─────────────────────────────────────────────────────────────
%  PALETA
% ─────────────────────────────────────────────────────────────
\definecolor{mainColor}{HTML}{00a499}
\definecolor{accentColor}{HTML}{EA7600}
\definecolor{bgLight}{HTML}{E6F7F6}
\definecolor{codeGray}{HTML}{2B2B2B}

% ─────────────────────────────────────────────────────────────
%  METROPOLIS
% ─────────────────────────────────────────────────────────────
\setbeamercolor{normal text}{fg=codeGray, bg=white}
\setbeamercolor{frametitle}{bg=mainColor, fg=white}
\setbeamercolor{alerted text}{fg=accentColor}
\setbeamercolor{progress bar}{fg=accentColor, bg=mainColor!25}
\setbeamercolor{block title}{bg=mainColor, fg=white}
\setbeamercolor{block body}{bg=bgLight, fg=codeGray}
\setbeamercolor{block title alerted}{bg=accentColor, fg=white}
\setbeamercolor{block body alerted}{bg=accentColor!10, fg=codeGray}
\metroset{progressbar=frametitle, sectionpage=none, numbering=fraction}

% ─────────────────────────────────────────────────────────────
%  IMAGEN CON FALLBACK AUTOMÁTICO
%  Uso: \img{image1.jpg}{Descripción}
%  → Si el archivo existe: muestra la imagen
%  → Si no existe:         muestra placeholder con el nombre del archivo
%  Para reemplazar: pon el archivo .jpg/.png junto al .tex
% ─────────────────────────────────────────────────────────────
\newcommand{\img}[3][0.95\linewidth]{%
  \begin{center}
  \IfFileExists{#2}{%
    \includegraphics[width=#1, keepaspectratio]{#2}%
  }{%
    \begin{tikzpicture}
      \node[draw=mainColor!60, fill=bgLight, rounded corners=4pt,
            minimum width=#1, minimum height=2.0cm, align=center,
            text=mainColor, font=\small\itshape]
        {\faImage\enspace\texttt{\detokenize{#2}}\\[3pt]
         \scriptsize\textcolor{codeGray}{#3}};
    \end{tikzpicture}%
  }%
  \end{center}
}

% ─────────────────────────────────────────────────────────────
%  CAJA DE CONCEPTO
% ─────────────────────────────────────────────────────────────
\newtcolorbox{conceptbox}[1]{
  colback=bgLight, colframe=mainColor,
  title={\small #1}, fonttitle=\bfseries,
  coltitle=white, colbacktitle=mainColor,
  arc=4pt, boxrule=0.6pt,
  left=4pt, right=4pt, top=2pt, bottom=2pt
}

\newcommand{\tech}[1]{\texttt{\textcolor{mainColor}{#1}}}

% ─────────────────────────────────────────────────────────────
%  PORTADA
% ─────────────────────────────────────────────────────────────
\title{Laboratorio Integrador\\de Geoinformática}
\subtitle{Semestre 1, 2026 · Universidad de Santiago de Chile}
\author{Byron Caices Lima \\ \small Ayudante --- Ingeniería Civil en Informática}
\date{\today}
% \titlegraphic{\includegraphics[height=1.2cm]{usach_logo.png}}

\begin{document}

\maketitle

% ─────────────────────────────────────────────────────────────
%  1. QUIÉN SOY
% ─────────────────────────────────────────────────────────────
\begin{frame}{Quién soy}
  \begin{columns}[T]
    \begin{column}{0.63\textwidth}
      \small
      \textbf{Formación}
      \begin{itemize}\setlength\itemsep{2pt}
        \item Estudiante de Ingeniería Civil en Informática, USACH (Licenciado).
        \item Me dedico principalmente a \textbf{Backend, Full-Stack e Infraestructura Cloud}.
        \item Mi base, igual que la de ustedes, es \textbf{informática}, no geografía.
              La diferencia es que ya cursé este laboratorio.
      \end{itemize}
      \vspace{0.25cm}
      \textbf{Mi proyecto en el ramo (como alumno)}
      \begin{itemize}\setlength\itemsep{2pt}
        \item \textit{Asesor inmobiliario virtual:} fusión de análisis geoespacial,
              datos censales y ML para hacer \textit{match} entre la vivienda ideal
              en Santiago y la realidad socioeconómica del usuario.
      \end{itemize}
    \end{column}
    \begin{column}{0.34\textwidth}
      % IMAGEN 1: Tu foto o avatar. Pon image1.jpg junto al .tex para reemplazar.
      \img{image1.jpg}{Foto o avatar del ayudante}
    \end{column}
  \end{columns}
\end{frame}

% ─────────────────────────────────────────────────────────────
%  2. ROL
% ─────────────────────────────────────────────────────────────
\begin{frame}{Mi rol: ¿en qué los voy a ayudar?}
  \begin{columns}[T]
    \begin{column}{0.48\textwidth}
      \begin{conceptbox}{\faCode\enspace Lo que \textit{sí} haré}
        \begin{itemize}\small\setlength\itemsep{1pt}
          \item Revisar estructura de código y repositorio
          \item Depurar errores de entorno y dependencias
          \item Guiar decisiones de arquitectura de software
          \item \textbf{Exigir buenas prácticas de ingeniería}
        \end{itemize}
      \end{conceptbox}
    \end{column}
    \begin{column}{0.48\textwidth}
      \begin{conceptbox}{\faGlobeAmericas\enspace Dónde el profe les puede ayudar +}
        \begin{itemize}\small\setlength\itemsep{1pt}
          \item Teoría de proyecciones cartográficas
          \item Interpretación de índices espectrales
          \item Metodología geoespacial profunda
          \item Validación del enfoque científico
        \end{itemize}
      \end{conceptbox}
    \end{column}
  \end{columns}
  \vspace{0.25cm}
  \begin{alertblock}{\faLightbulb\enspace Idea simple}
    \centering\small
    Duda de \textbf{código, entorno o estructura} \textrightarrow{} yo.\quad
    Duda de \textbf{geografía o metodología} \textrightarrow{} el profesor.
  \end{alertblock}
\end{frame}

% ─────────────────────────────────────────────────────────────
%  3. TIPOS DE DATOS
% ─────────────────────────────────────────────────────────────
\begin{frame}{Geoinformática para Informáticos: Tipos de Datos}
  \begin{columns}[T]
    \begin{column}{0.48\textwidth}
      \begin{conceptbox}{\faDrawPolygon\enspace Datos Vectoriales}
        \small
        \textbf{¿Qué son?} Geometrías discretas: puntos, líneas y polígonos.\\[3pt]
        \textbf{Analogía IT:} Filas en SQL con columna \tech{geometry}.\\[3pt]
        \textbf{Librerías:} \tech{GeoPandas} · \tech{Shapely}
      \end{conceptbox}
      \vspace{0.1cm}
      % IMAGEN 2: Mapa vectorial. Pon image2.jpg junto al .tex para reemplazar.
      \img[0.88\linewidth]{image2.jpeg}{Mapa vectorial: zonas censales de Santiago}
    \end{column}
    \begin{column}{0.48\textwidth}
      \begin{conceptbox}{\faTh\enspace Datos Raster}
        \small
        \textbf{¿Qué son?} Grillas espaciales: imágenes satelitales, temperatura, NDVI.\\[3pt]
        \textbf{Analogía IT:} Array NumPy 2D/3D donde cada celda tiene coordenada.\\[3pt]
        \textbf{Librerías:} \tech{Rasterio} · \tech{Xarray}
      \end{conceptbox}
      \vspace{0.1cm}
      % IMAGEN 3: Imagen Sentinel. Pon image3.jpg junto al .tex para reemplazar.
      \img[0.88\linewidth]{image3.jpg}{Imagen Sentinel-2 RGB de Santiago}
    \end{column}
  \end{columns}
\end{frame}

% ─────────────────────────────────────────────────────────────
%  4. EPSG:32719
% ─────────────────────────────────────────────────────────────
\begin{frame}{Sistema de Referencia Espacial: ¿Qué es EPSG:32719?}
  \begin{columns}[T]
    \begin{column}{0.57\textwidth}
      \small
      La Tierra es esférica. Un mapa es plano. \textbf{Proyectar} es la función
      que convierte una en el otro, y cada proyección tiene un código
      \textbf{EPSG} estandarizado.\\[6pt]
      \begin{itemize}\setlength\itemsep{4pt}
        \item \textbf{EPSG:4326 --- WGS84}\\
              Coordenadas en \textit{grados} (lat / lon).
              Para \textbf{almacenar y compartir}.
              Es el sistema de tu GPS y Google Maps.
        \item \textbf{EPSG:32719 --- UTM Zona 19S}\\
              Coordenadas en \textit{metros}.
              Para \textbf{calcular distancias, áreas y buffers}
              sobre Chile continental.
      \end{itemize}
      \vspace{0.2cm}
      \begin{alertblock}{\faSyncAlt\enspace Regla práctica}
        \small
        Guarda en \tech{4326}.\quad Calcula en \tech{32719}.\\[2pt]
        \tech{gdf = gdf.to\_crs(epsg=32719)}
      \end{alertblock}
    \end{column}
    \begin{column}{0.40\textwidth}
      % IMAGEN 4: Diagrama UTM. Pon image4.jpg junto al .tex para reemplazar.
      \img{image4.jpg}{Diagrama UTM Zona 19S sobre Chile}
    \end{column}
  \end{columns}
\end{frame}

% ─────────────────────────────────────────────────────────────
%  5. ENTREGABLES
% ─────────────────────────────────────────────────────────────
\begin{frame}{Reglas del Juego: Entregables del Proyecto}
  \begin{columns}[T]
    \begin{column}{0.57\textwidth}
      \small
      Todo proyecto debe cumplir con los siguientes requisitos técnicos:
      \vspace{0.15cm}
      \begin{itemize}\setlength\itemsep{4pt}
        \item[\textcolor{accentColor}{\faCompass}]
          Proyección estándar \textbf{EPSG:32719} para todos los cálculos en Chile.
        \item[\textcolor{accentColor}{\faMap}]
          Mínimo \textbf{5 mapas temáticos} (título, leyenda, escala, norte)
          y \textbf{8 gráficos}.
        \item[\textcolor{accentColor}{\faFileAlt}]
          Informe técnico en \LaTeX{} o Markdown: \textbf{25--35 páginas}.
        \item[\textcolor{accentColor}{\faChartBar}]
          \textbf{Dashboard interactivo} con Folium, Streamlit o Panel.
      \end{itemize}
    \end{column}
    \begin{column}{0.40\textwidth}
      % IMAGEN 5: Mapa temático. Pon image5.jpg junto al .tex para reemplazar.
      \img{image5.jpg}{Ejemplo de mapa temático con leyenda y escala}
    \end{column}
  \end{columns}
\end{frame}

% ─────────────────────────────────────────────────────────────
%  6. PIPELINE REPRODUCIBLE
% ─────────────────────────────────────────────────────────────
\begin{frame}{El Estándar: Pipeline 100\,\% Reproducible}
  \begin{columns}[T]
    \begin{column}{0.67\textwidth}
      \begin{alertblock}{\faExclamationTriangle\enspace Regla de oro}
        \small\textbf{No instalaremos nada manualmente} para evaluar su proyecto.
        Todo debe correr desde cero con un solo comando.
      \end{alertblock}
      \vspace{0.2cm}
      \small El repositorio debe incluir:
      \begin{enumerate}\setlength\itemsep{2pt}
        \item \textbf{GitHub público} con historial de commits real.
        \item \textbf{Docker + Docker Compose} para aislar el entorno
              y resolver conflictos de dependencias geográficas.
        \item Script \tech{setup.sh} / \tech{run.bat} que:
              \begin{itemize}\setlength\itemsep{1pt}
                \item Levante los contenedores.
                \item Instale dependencias geográficas.
                \item Levante la BD y ejecute el autopoblado.
              \end{itemize}
      \end{enumerate}
    \end{column}
    \begin{column}{0.4\textwidth}
      % IMAGEN 6: Docker Compose. Pon image6.jpg junto al .tex para reemplazar.
      \img{image6.jpg}{Diagrama Docker Compose: Jupyter + PostGIS + Streamlit}
    \end{column}
  \end{columns}
\end{frame}

% ─────────────────────────────────────────────────────────────
%  7a. ARQUITECTURA REPO — 1 / 2
% ─────────────────────────────────────────────────────────────
\begin{frame}[fragile]{Arquitectura del Repositorio (1 / 2)}
\footnotesize
\begin{verbatim}
laboratorio_integrador/
├── README.md                  # Descripción del proyecto
├── requirements.txt           # Dependencias Python
├── docker-compose.yml         # Orquestación de contenedores
├── .env                       # Variables de entorno  ← NO subir a Git
├── .gitignore                 # Excluye datos pesados (rásteres, etc.)
│
├── docker/
│   ├── jupyter/               # Imagen personalizada de Jupyter
│   ├── postgis/               # Inicialización de PostGIS
│   └── web/                   # Imagen de la app web
│
├── data/
│   ├── raw/                   # Datos originales sin procesar
│   ├── processed/             # Datos procesados y limpios
│   └── external/              # Fuentes externas (INE, Copernicus…)
│
└── notebooks/
    ├── 01_data_acquisition.ipynb
    ├── 02_exploratory_analysis.ipynb
    ├── 03_geostatistics.ipynb
    ├── 04_machine_learning.ipynb
    └── 05_results_synthesis.ipynb
\end{verbatim}
\end{frame}

% ─────────────────────────────────────────────────────────────
%  7b. ARQUITECTURA REPO — 2 / 2
% ─────────────────────────────────────────────────────────────
\begin{frame}[fragile]{Arquitectura del Repositorio (2 / 2)}
\footnotesize
\begin{verbatim}
laboratorio_integrador/                         (continuación)
│
├── scripts/
│   ├── download_data.py       # Descarga automatizada de datos
│   ├── process_data.py        # Pipeline de procesamiento
│   ├── spatial_analysis.py    # Análisis espacial reutilizable
│   └── utils.py               # Funciones auxiliares comunes
│
├── app/
│   ├── main.py                # Punto de entrada Streamlit
│   ├── pages/                 # Vistas del dashboard
│   └── components/            # Componentes reutilizables
│
├── outputs/
│   ├── figures/               # Gráficos y mapas generados
│   ├── models/                # Modelos ML serializados
│   └── reports/               # Informes y documentos finales
│
└── docs/
    ├── guia_usuario.md
    ├── arquitectura.md
    └── api_reference.md
\end{verbatim}
\end{frame}

% ─────────────────────────────────────────────────────────────
%  8. CALENDARIO
% ─────────────────────────────────────────────────────────────
\begin{frame}{Calendario de Hitos Clave}
  \vspace{-0.1cm}
  \begin{center}
  \begin{tikzpicture}[font=\small, >=Stealth]
    \draw[->, line width=1.2pt, color=mainColor] (0,0) -- (12.4,0);
    \foreach \x/\fecha/\hitoa/\hitob in {
      0.3/{14 abr}/{Grupos y}/{3 desafíos},
      2.8/{28 abr}/{Desafío}/{final},
      5.5/{26 may}/{PEP 1}/{ },
      8.5/{25 jun}/{Pipeline}/{en Docker},
      11.8/{23 jul}/{PEP 2}/{Demo Day}%
    }{
      \fill[accentColor] (\x,0) circle (3.5pt);
      \node[above=5pt, align=center, text=codeGray, font=\scriptsize]
        at (\x,0) {\hitoa\\\hitob};
      \node[below=5pt, align=center, text=accentColor, font=\scriptsize\bfseries]
        at (\x,0) {\fecha};
    }
  \end{tikzpicture}
  \end{center}
  \vspace{0.15cm}
  \begin{columns}
    \begin{column}{0.48\textwidth}
      \begin{block}{\faClipboardList\enspace PEP 1 --- 26 de mayo}
        \small Documento de validación + Avance pipeline + inicio de EDA.
      \end{block}
    \end{column}
    \begin{column}{0.48\textwidth}
      \begin{block}{\faTrophy\enspace PEP 2 --- 23 de julio}
        \small Dashboard final completo + defensa oral (Demo Day).
      \end{block}
    \end{column}
  \end{columns}
\end{frame}

% ─────────────────────────────────────────────────────────────
%  9. TAREAS PARA HOY
% ─────────────────────────────────────────────────────────────
% ─────────────────────────────────────────────────────────────
%  9. TAREAS PARA HOY
% ─────────────────────────────────────────────────────────────
\begin{frame}{Para Hoy y la Próxima Clase}
  \begin{columns}[T]
    \begin{column}{0.67\textwidth}
      \small
      \begin{enumerate}\setlength\itemsep{4pt}
        \item \textbf{Armar los grupos} (1 a 3 integrantes).
        \item \textbf{Revisar el catálogo:} analicen los 12 desafíos
              disponibles. Elijan según lo que prefieran trabajar.
        \item \textbf{Crear el repositorio} en GitHub con la
              estructura de los slides anteriores.
        \item \textbf{Instalar Docker Desktop} en sus máquinas
              y verificar que corre.
        \item \textbf{Subir el link del repositorio} a la sección
              de Laboratorio habilitada en \textbf{UVirtual}.
      \end{enumerate}
      \vspace{0.15cm}
      \begin{alertblock}{}
        \centering\small
        Al terminar la sesión: cada grupo debe tener
        \textbf{repo creado, al menos un commit y el link entregado en UVirtual}.
      \end{alertblock}
    \end{column}
    \begin{column}{0.40\textwidth}
      % IMAGEN 7: Repo GitHub. Pon image7.jpg junto al .tex para reemplazar.
      \img{image7.jpg}{Captura GitHub: repo con estructura inicial}
      \vspace{0.2cm}
      \centering\large\textbf{¡Manos a la obra!}
    \end{column}
  \end{columns}
\end{frame}

\end{document}